% Weekly exercise sheet for Lecture 12; also included in the synthesis sheet.
\ifdefined\ExercisesSynthesis
\else
\documentclass[11pt]{article}
\usepackage[utf8]{inputenc}
\usepackage[T1]{fontenc}
\usepackage[english]{babel}
\usepackage{lmodern,microtype}
\usepackage[margin=0.9in]{geometry}
\usepackage{amsmath,amssymb,amsthm,mathtools}
\usepackage{enumitem,hyperref}
\newcommand{\E}{\mathbb E}
\newcommand{\Qbb}{\mathbb Q}
\newcommand{\dd}{\,\mathrm d}
\newcommand{\F}{\mathcal F}
\theoremstyle{definition}
\newtheorem{exercise}{Exercise}[section]
\newlist{exparts}{enumerate}{1}
\setlist[exparts,1]{label=\textbf{(\alph*)},itemsep=0.3em,topsep=0.4em,leftmargin=2em}
\newcommand{\lecturesummary}[1]{\begin{quote}\small\textbf{Focus.} #1\end{quote}}
\title{Stochastic Calculus and Option Pricing\\Lecture 12 Exercises}
\author{}
\date{}
\begin{document}
\maketitle
\vspace{-1.2em}
\setcounter{section}{11}
\fi

\section{Week 12 --- Beyond Black--Scholes}
\lecturesummary{The questions isolate covariance in random discounting and
the risk-spanning consequence of adding a volatility factor.}

\begin{exercise}[Random discounting and conditional covariance]
Under a risk-neutral measure, let
\[
D_{t,T}=\exp\left(-\int_t^T r_u\,\dd u\right)
\]
and let a claim pay \(H_T\geq0\) at \(T\).
\begin{exparts}
  \item Write its time-\(t\) value as a conditional expectation and
  decompose it into a product of conditional means plus a conditional
  covariance term.
  \item Give a sufficient condition under which replacing the random
  discount factor by its conditional mean times the expected payoff is
  valid.
  \item Explain why a positive covariance between discounting and payoff
  raises the value relative to that product approximation.
\end{exparts}
\end{exercise}

\begin{exercise}[A second Brownian shock adds a hedge direction]
Suppose \(V_t>0\), and the stock and variance satisfy
\[
\frac{\dd S_t}{S_t}=r\dd t+\sqrt{V_t}\dd W_t^S,\qquad
\dd V_t=b(t,V_t)\dd t+\eta(t,V_t)\dd W_t^V,
\]
with \(|\rho|<1\) between \(W^S\) and \(W^V\), and \(\eta\ne0\).
\begin{exparts}
  \item Explain why the two Brownian shocks span two independent local
  risk directions.
  \item A desk trades only the stock and cash. Which risk direction can it
  hedge instantaneously, and which remains?
  \item What feature would a second traded instrument need in order to
  complete the local hedge? State the condition in terms of its exposure
  vector, rather than naming a particular product.
\end{exparts}
\end{exercise}

\ifdefined\ExercisesSynthesis\else
\end{document}
\fi
